\documentclass[lettersize,journal]{IEEEtran}

\usepackage{amsmath,amsfonts}
\usepackage{algorithm}
\usepackage{algorithmic}
\usepackage{array}
\usepackage{booktabs}
\usepackage{textcomp}
\usepackage{stfloats}
\usepackage{url}
\usepackage{graphicx}
\usepackage{cite}

\graphicspath{{./figures/}}
\newcommand{\maybeincludegraphics}[2][]{\includegraphics[#1]{#2}}

\hyphenation{op-tical net-works semi-conduc-tor IEEE-Xplore}

\begin{document}

\title{Trigger-Free Jailbreak Backdoor Mitigation for Large Language Models via Confidence-Gated Perturbation-Sensitive Pruning}

\author{Anonymous Authors}

\markboth{IEEE Transactions Draft}%
{Anonymous Authors: Trigger-Free Jailbreak Backdoor Mitigation for LLMs}

\maketitle

\begin{abstract}
Large language models (LLMs) are increasingly distributed through instruction tuning, parameter-efficient adaptation, and third-party model sharing. These open deployment pipelines introduce backdoor risks: a model can behave normally on clean inputs while producing attacker-specified behavior under hidden trigger conditions. Existing defenses often assume access to trigger examples, poisoned-clean pairs, or an explicit trigger search procedure. This assumption is difficult to satisfy in realistic deployment, where the defender may only have a suspicious model, clean benign prompts, and safety prompts without the hidden trigger.

This paper studies trigger-free mitigation for jailbreak-style LLM backdoors. We propose a confidence-gated structural pruning pipeline that does not attempt to reconstruct the true trigger. Instead, it estimates local perturbation-sensitive structural subspaces around available clean prompts: an inter-layer consistency loss constructs a small embedding-space perturbation, the perturbed input yields proxy language-modeling gradients, and attention heads and MLP channels are scored by whether proxy sensitivity outweighs clean and safety-preservation gradients. Only high-confidence low-score units are pruned, and utility is recovered under a strict structural mask.

Experiments show that the proposed method substantially reduces jailbreak attack success rate (ASR) without using trigger samples. On BEAT Llama-3.1 jailbreak backdoors, the method reduces word-trigger ASR from 0.9250 to 0.1417, phrase-trigger ASR to 0.0917, and long-trigger ASR to 0.1833 under a fixed runtime environment. A tuned CROW direct-repair baseline reduces ASR only modestly, while auditable matched-random controls show that the scored pruning plan outperforms layer/component-matched and score-band-matched random pruning under the same recovery protocol. Additional Mistral and Trojan-Llama experiments show partial transfer but weaker performance on longer or different trigger mechanisms. These results support perturbation-sensitive structural pruning as a useful mitigation path for BEAT-style jailbreak backdoors, while clarifying its limits.
\end{abstract}

\begin{IEEEkeywords}
Large language models, backdoor defense, trigger-free mitigation, structured pruning, adversarial consistency, model safety.
\end{IEEEkeywords}

\section{Introduction}
\IEEEPARstart{L}{arge} language models (LLMs) are becoming core components of interactive systems and automated workflows. Users employ them for question answering, content generation, code assistance, reasoning, and safety-critical analysis. As model capabilities improve, behavioral reliability becomes increasingly important: a single jailbreak on a harmful request can break a safety boundary, and a trigger-activated behavior can create systematic deployment risk.

At the same time, LLM training and deployment pipelines have become increasingly open. Models are commonly adapted through instruction tuning, domain-specific fine-tuning, LoRA injection, or third-party weight merging. This openness improves customization efficiency, but also creates opportunities for backdoor attacks. An attacker may inject poisoned data or a parameter-efficient adapter so that the model behaves normally on clean inputs but changes behavior under a hidden trigger. Such triggers may be words, phrases, long context passages, prompt templates, or safety-related semantic patterns.

Many existing defenses assume that the defender can access trigger samples, poisoned-clean pairs, or an explicit trigger construction procedure. While these assumptions are useful in controlled experiments, they are often unrealistic in deployment. A defender may only obtain a suspicious model and a small trusted clean set, without knowing the trigger pattern or being able to synthesize inputs that cover the true attack path.

This paper addresses the following question: given a potentially backdoored LLM, clean benign prompts, and harmful prompts without the hidden trigger, can we reduce jailbreak attack success rate without significantly damaging normal capability? This setting is challenging for three reasons. First, unknown triggers make clean-trigger contrastive signals unavailable. Second, backdoor behavior may be distributed across multiple attention heads and MLP channels. Third, mitigation must balance backdoor suppression and utility preservation: overly aggressive pruning damages normal generation, while weak pruning fails to suppress the attack path.

We propose \emph{Trigger-Free Pruning Defense}, a structured mitigation framework for unknown-trigger scenarios. The key idea is that although the real trigger is unavailable, backdoor paths often correspond to unstable internal representation directions. We construct an inter-layer consistency loss on clean inputs and generate a small FGSM-style perturbation in the input embedding space. The perturbed input is not a trigger, but it amplifies units that are sensitive to abnormal representation shifts. We then use the clean gradient and the perturbation-proxy gradient to score structural units, prune suspicious attention heads and MLP channels, and recover model utility under a strict mask.

The contributions of this paper are as follows.
\begin{itemize}
    \item We formulate a trigger-free structured mitigation pipeline for BEAT-style jailbreak backdoors that combines perturbation-sensitive scoring, confidence-gated pruning, and strict-mask recovery fine-tuning.
    \item We adapt clean-data adversarial consistency perturbations into a proxy-gradient scoring module for head/channel-level pruning. The proxy is not assumed to reconstruct the true trigger; it estimates local vulnerable subspaces that may overlap with trigger-induced backdoor pathways.
    \item We introduce a head/channel-level score with clean-gradient protection, optional harmful-no-trigger safety protection, and an absolute cosine proxy penalty, together with a high-confidence score cap that avoids large-budget low-confidence pruning.
    \item We compare against a tuned CROW-style direct perturbation-repair baseline, showing that direct consistency training gives only weak ASR reduction on BEAT jailbreak backdoors and motivating the conversion of perturbation sensitivity into structural pruning decisions.
    \item We provide auditable matched-random controls showing that the scored plan outperforms uniform, layer/component-matched, and score-band-matched random pruning under the same recovery protocol.
\end{itemize}

\section{Related Work}

\subsection{Backdoor Attacks on Language Models}
Backdoor attacks aim to make a model behave normally on most inputs while producing attacker-desired behavior under specific trigger conditions~\cite{badnets,trojaning,hidden_trigger,beat}. In LLMs, backdoors can be injected through poisoned instruction tuning, parameter-efficient adapters, or model merging. Unlike classification backdoors, generative-model backdoors can manifest as fixed target text, jailbreak behavior, hidden bias, or task-specific generation patterns.

LLM backdoors are strongly context-dependent. Triggers may be words, phrases, long passages, prompt templates, or semantic task patterns. After activation, the model may produce attacker-specified content or bypass safety refusal on harmful requests. Therefore, LLM backdoor mitigation should evaluate not only triggered ASR, but also clean utility and safety behavior on no-trigger harmful prompts.
This context dependence is related to broader evidence that prompting and adversarial suffixes can expose brittle safety boundaries in language models~\cite{universal_triggers,autoprompt,zou2023universal,gptfuzzer,jailbroken}, and that some LLM backdoors can persist through ordinary safety training~\cite{sleeper_agents}.

\subsection{Backdoor Defenses}
Existing defenses include fine-tuning-based mitigation, trigger inversion, activation-based detection, pruning, and robustness regularization~\cite{neuralcleanse,spectral,activation_clustering,finepruning,crow}. Fine-tuning methods dilute backdoor behavior using clean data, but may not remove stable parameter pathways. Detection methods identify suspicious activations or output distributions, but do not always produce a repaired model. Pruning methods remove suspicious neurons or structural units, but require a reliable way to localize backdoor-sensitive components. Robustness and consistency methods regularize behavior under perturbations, but may target general instability rather than backdoor-specific pathways~\cite{fgsm,pgd,adv_training_free}.

Our structural view is also connected to neural-network pruning and sparsity work, where weights, attention heads, or larger structured blocks can be removed while preserving much of the original function~\cite{magnitude_pruning,lottery,head_pruning,movement_pruning,llm_pruner}. In contrast to generic compression, however, our goal is not to minimize model size; it is to remove high-confidence perturbation-sensitive structures while preserving clean and no-trigger safety behavior.

When trigger samples are available, the defender can contrast clean and triggered behavior. In trigger-free settings, this contrast is unavailable. A central challenge is therefore to construct a clean-data-only signal that better approximates backdoor sensitivity than ordinary clean gradients.

\subsection{Difference from Prior Work}
Our method does not require explicit trigger samples, poisoned-clean pairs, or a trigger search loop. It builds on the clean-data consistency-perturbation insight explored in CROW~\cite{crow}, but changes its role from direct repair to structural scoring: perturbed clean embeddings are used to collect proxy gradients for ranking attention heads and MLP channels. We do not claim that this proxy recovers the hidden trigger. Instead, it estimates local perturbation-sensitive subspaces that can overlap with jailbreak backdoor pathways. Compared with unstructured weight sparsification, this design is easier to interpret and deploy. Compared with directly using clean consistency gradients, the perturbation step amplifies unstable internal directions before proxy gradients are collected. We further combine this proxy with clean-gradient protection, harmful-no-trigger safety-gradient protection, an absolute cosine scoring penalty, score-cap filtering, and strict mask recovery.

\section{Method}

\subsection{Problem Definition}
Let $M_\theta$ be a suspicious LLM. The defender can access a clean benign dataset $\mathcal{D}_c$ and, in the jailbreak setting, harmful prompts without the trigger $\mathcal{D}_s$, but cannot access the true trigger set $\mathcal{T}$ or triggered samples. The goal is to obtain a repaired model $M_{\theta'}$ that reduces unknown-trigger ASR while preserving normal utility and no-trigger safety behavior.

We denote model structural units by $u \in \mathcal{U}$. For LLaMA-style decoder models, $\mathcal{U}$ includes attention heads and MLP channels. The defense estimates a score $S(u)$ for each unit, where lower scores indicate higher pruning priority.

\subsection{Adversarial Consistency Proxy}
Following the clean-data perturbation intuition in prior consistency-based repair work~\cite{crow} and the broader use of adversarial input perturbations~\cite{fgsm,pgd}, we use perturbations as a diagnostic proxy rather than as real trigger samples.
Given a clean input $x$, the model produces hidden states $h_\ell$ at layer $\ell$. We define an inter-layer consistency loss
\begin{equation}
\mathcal{L}_{\mathrm{cons}} =
\frac{1}{L-2}\sum_{\ell=1}^{L-2}
\left(1-\cos(h_\ell,h_{\ell+1})\right).
\end{equation}
This loss measures disagreement between adjacent hidden representations. In implementation, we average over adjacent internal hidden states while excluding the embedding state and the final output state, so the summation upper bound is $L-2$. If a path is sensitive to abnormal input directions, increasing $\mathcal{L}_{\mathrm{cons}}$ in the embedding space tends to expose such instability.

We generate an FGSM-style perturbation in the input embedding space:
\begin{equation}
\delta =
\epsilon_p \cdot
\mathrm{sign}\left(\nabla_{E(x)} \mathcal{L}_{\mathrm{cons}}\right),
\end{equation}
\begin{equation}
E(x)^{\mathrm{adv}} = E(x)+\delta,
\end{equation}
where $E(x)$ denotes token embeddings and $\epsilon_p$ is the perturbation strength. The perturbed input is not treated as a trigger; it is a proxy for probing unstable internal directions.

The proxy gradient is then collected from a second language-modeling forward pass on $E(x)^{\mathrm{adv}}$, as illustrated in Fig.~\ref{fig:proxy_gradient_principle}. Concretely, the implementation sets the token labels equal to the input token ids and relies on the standard causal-LM shift: logits at each position predict the next token, and the language-modeling cross-entropy is computed against the shifted token labels. Back-propagating this perturbed-input loss gives the parameter gradients used as $g_p(u)$. Thus, the proxy gradient is not the gradient of $\mathcal{L}_{\mathrm{cons}}$ itself. The consistency loss only determines the embedding perturbation direction; the structural proxy score is derived from the language-modeling loss after applying that perturbation.

For selected layers $\mathcal{S}$, we further use a hidden-state alignment loss:
\begin{equation}
\mathcal{L}_{\mathrm{align}} =
\sum_{\ell \in \mathcal{S}} \lambda_\ell
\left(1-\cos(H_\ell^c,H_\ell^{\mathrm{adv}})\right),
\end{equation}
where $H_\ell^c$ and $H_\ell^{\mathrm{adv}}$ are clean and perturbed hidden states. This alignment objective is used during optional pre-alignment and recovery to stabilize clean-neighborhood representation behavior. To avoid confusion, we distinguish three related quantities: $\mathcal{L}_{\mathrm{cons}}$ is used only to generate the embedding perturbation direction; $\mathcal{L}_{\mathrm{align}}$ is the clean-versus-perturbed hidden-state alignment loss used during recovery, with per-layer weights $\lambda_\ell$; and the scalar $\lambda_{\mathrm{align}}$ in the recovery objective is the overall weight of $\mathcal{L}_{\mathrm{align}}$ in the recovery loss. For structural scoring, the implementation uses $\mathcal{L}_{\mathrm{cons}}$ to construct $E(x)^{\mathrm{adv}}$, then obtains the proxy gradient from the language-modeling loss of the perturbed input rather than directly from $\mathcal{L}_{\mathrm{align}}$.

\begin{figure*}[!t]
\centering
\maybeincludegraphics[width=0.94\textwidth]{fig_proxy_gradient_principle.pdf}
\caption{Perturbation proxy gradient construction. A clean input first produces hidden states whose inter-layer similarity loss yields an embedding-gradient direction. The perturbed embedding is then fed through the same model, and its next-token language-modeling cross-entropy against shifted token labels is back-propagated to obtain proxy gradients. The perturbation is a trigger-free diagnostic signal rather than a recovered textual trigger.}
\label{fig:proxy_gradient_principle}
\end{figure*}

\subsection{Structural Scoring}
For each structural unit $u$, we collect three gradients. The clean gradient $g_c(u)$ reflects normal-task contribution, the proxy gradient $g_p(u)$ reflects sensitivity to the perturbation proxy, and, when harmful-no-trigger prompts are available, the safety gradient $g_s(u)$ protects no-trigger refusal behavior. Fig.~\ref{fig:clean_safe_gradient_principle} shows how the clean and safety losses are obtained. Both are ordinary causal language-modeling losses: a benign prompt or a harmful-no-trigger prompt is forwarded through the model, next-token logits are compared with shifted token labels, and the loss is back-propagated. The difference between $g_c$ and $g_s$ is therefore the prompt distribution, not the loss type. The supervised refusal target used during recovery is separate from this scoring-stage safety gradient.

\begin{figure*}[!t]
\centering
\maybeincludegraphics[width=0.94\textwidth]{fig_clean_safe_gradient_principle.pdf}
\caption{Clean and safety gradient collection for structural scoring. For clean prompts and harmful prompts without triggers, the model computes next-token language-modeling logits and cross-entropy against shifted token labels. Back-propagating the clean loss yields $g_c$, while back-propagating the harmful-no-trigger loss yields $g_s$. These gradients are later aggregated over attention heads and MLP channels to form $G_c$ and $G_s$.}
\label{fig:clean_safe_gradient_principle}
\end{figure*}

The gradients are converted into structural magnitudes by grouping parameter gradients according to the unit definition. For an attention head, we flatten and aggregate the gradients of its corresponding query, key, value, and output projection slices. For an MLP channel, we aggregate the gate, up, and down projection gradients. The unit-level magnitude is the mean absolute value of the flattened gradient vector. Let
\begin{equation}
\begin{aligned}
G_c(u) &= \mathrm{mag}(g_c(u)),\\
G_p(u) &= \mathrm{mag}(g_p(u)),\\
G_s(u) &= \mathrm{mag}(g_s(u)).
\end{aligned}
\end{equation}
The directional cosine is
\begin{equation}
\cos\theta(u)=
\frac{\langle g_c(u),g_p(u)\rangle}
{\|g_c(u)\|_2\|g_p(u)\|_2+\xi}.
\end{equation}
The main paper-facing score is
\begin{equation}
S(u)=\alpha\left(G_c(u)+\alpha_{\mathrm{safe}}G_s(u)\right)
-\beta |G_p(u)\cos\theta(u)|.
\end{equation}
The first term protects units important for clean behavior and, when enabled, units important for no-trigger harmful-prompt refusal. The second term lowers the score for units strongly responding to the perturbation proxy. The absolute value is important: if clean and proxy gradients point in opposite directions, the unit should not be rewarded, because such a direction indicates deviation from normal behavior. In implementation, $\mathrm{mag}(\cdot)$ is the mean absolute gradient magnitude over the flattened parameters belonging to a structural unit, while $\cos\theta$ is computed from the flattened clean and proxy gradient vectors. Setting $\alpha_{\mathrm{safe}}=0$ recovers the clean/proxy-only score.

The implementation also supports a harmful-context proxy term,
\begin{equation}
S_{\mathrm{harm}}(u)=S(u)-\beta_h |G_{p,h}(u)\cos\theta_h(u)|,
\end{equation}
where $G_{p,h}$ is computed by applying the same perturbation-proxy procedure to harmful-no-trigger prompts. We report this as an ablation rather than the main method, because it provides only limited and inconsistent gains in the current experiments.

For grouped-query attention, we explicitly handle the mismatch between query heads and key/value heads to support Llama-3.1-style architectures.

\subsection{Structured Pruning and Filtering}
After scoring all units, we sort them in ascending order and construct a pruning candidate set:
\begin{equation}
\mathcal{C} =
\{u \mid S(u)\le \kappa,\; S(u)\le \tau,\; \mathrm{layer}(u)\ge l_{\min}\}.
\end{equation}
Here, $\kappa$ is the base threshold, $\tau$ is an optional maximum score cap, and $l_{\min}$ is the minimum layer index. We then select the first $B$ units from $\mathcal{C}$. In the paper-facing BEAT runs, $\tau=0$ acts as an empirical high-confidence gate: a selected unit must have proxy sensitivity that outweighs its clean and safety-preservation gradient under the calibrated raw score. This should not be interpreted as a universal theoretical boundary. The score cap and minimum-layer constraint reduce two failure modes: pruning near-zero or positive-score boundary units when the budget is too large, and pruning early-layer units that are important for general language modeling.

Pruning is implemented with structured masks, following the motivation that head-level and channel-level structures are interpretable pruning targets in transformer models~\cite{head_pruning,movement_pruning,llm_pruner}. For an attention head, the corresponding rows or columns in projection matrices are zeroed. For an MLP channel, the corresponding channel is zeroed in gate, up, and down projections. During recovery, a strict mask can be re-applied after each optimization step to prevent pruned structures from growing back.

\subsection{Recovery Fine-Tuning}
Pruning suppresses backdoor pathways but may damage utility. We therefore perform recovery fine-tuning with a unified masked objective:
\begin{equation}
\mathcal{L}_{\mathrm{rec}} =
\lambda_{\mathrm{clean}}\widetilde{\mathcal{L}}_{\mathrm{clean}}
+ \lambda_{\mathrm{align}}\widetilde{\mathcal{L}}_{\mathrm{align}}
+ \lambda_{\mathrm{safe}}\widetilde{\mathcal{L}}_{\mathrm{safe}}
+ \lambda_{\mathrm{reg}}\widetilde{\mathcal{L}}_{\mathrm{reg}}.
\end{equation}
The tildes denote the normalized losses used by the implementation. $\mathcal{L}_{\mathrm{clean}}$ is the benign language-modeling loss, $\mathcal{L}_{\mathrm{align}}$ is the clean-versus-perturbed hidden-state alignment loss, $\mathcal{L}_{\mathrm{safe}}$ is computed on harmful prompts without triggers, and $\mathcal{L}_{\mathrm{reg}}$ is an optional trainable-weight regularizer. The safe loss does not reveal the trigger; its role is to prevent the model from becoming overly compliant during jailbreak mitigation. When no safety-preservation split is used, $\lambda_{\mathrm{safe}}=0$.

\begin{algorithm}[!t]
\caption{Trigger-Free Confidence-Gated Pruning}
\label{alg:tfpd}
\begin{algorithmic}[1]
\REQUIRE Suspicious model $M_\theta$; clean set $\mathcal{D}_c$; harmful-no-trigger set $\mathcal{D}_s$; perturbation strength $\epsilon_p$; score weights $\alpha,\beta,\alpha_{\mathrm{safe}}$; threshold $\kappa$, score cap $\tau$, min layer $l_{\min}$, budget $B$
\ENSURE Repaired model $M_{\theta'}$
\STATE \textit{// Stage 1: perturbation-proxy scoring}
\FOR{each clean input $x \in \mathcal{D}_c$}
  \STATE Compute $\mathcal{L}_{\mathrm{cons}}$ from hidden states; $\delta \gets \epsilon_p\,\mathrm{sign}(\nabla_{E(x)}\mathcal{L}_{\mathrm{cons}})$
  \STATE Forward $E(x)^{\mathrm{adv}}=E(x)+\delta$; back-prop LM loss to accumulate proxy gradient $g_p$
\ENDFOR
\STATE Accumulate clean gradient $g_c$ on $\mathcal{D}_c$ and safety gradient $g_s$ on $\mathcal{D}_s$
\FOR{each structural unit $u \in \mathcal{U}$}
  \STATE $G_c,G_p,G_s \gets \mathrm{mag}(g_c),\mathrm{mag}(g_p),\mathrm{mag}(g_s)$; compute $\cos\theta(u)$
  \STATE $S(u) \gets \alpha\!\left(G_c+\alpha_{\mathrm{safe}}G_s\right) - \beta\,|G_p\cos\theta(u)|$
\ENDFOR
\STATE \textit{// Stage 2: confidence-gated structured pruning}
\STATE $\mathcal{C} \gets \{u: S(u)\le\kappa,\ S(u)\le\tau,\ \mathrm{layer}(u)\ge l_{\min}\}$
\STATE Select first $B$ units of $\mathcal{C}$ in ascending $S$; build structural mask; zero pruned units
\STATE \textit{// Stage 3: masked recovery}
\FOR{each recovery step}
  \STATE Update $\theta$ on $\mathcal{L}_{\mathrm{rec}}=\lambda_{\mathrm{clean}}\widetilde{\mathcal{L}}_{\mathrm{clean}}+\lambda_{\mathrm{align}}\widetilde{\mathcal{L}}_{\mathrm{align}}+\lambda_{\mathrm{safe}}\widetilde{\mathcal{L}}_{\mathrm{safe}}+\lambda_{\mathrm{reg}}\widetilde{\mathcal{L}}_{\mathrm{reg}}$
  \STATE Re-apply strict mask to keep pruned units at zero
\ENDFOR
\RETURN $M_{\theta'}$
\end{algorithmic}
\end{algorithm}

\section{Experimental Setup}

\subsection{Models and Backdoor Types}
The main experiments evaluate BEAT Llama-3.1 jailbreak backdoors~\cite{beat}, where the triggered model bypasses safety refusal on harmful prompts. We also evaluate Mistral-3-7B BEAT variants and Trojan-Llama-8B as stress tests for cross-model and cross-mechanism transfer.

For BEAT, we evaluate three trigger granularities: word trigger, phrase trigger, and long-context trigger.

\subsection{Metrics}
For jailbreak-style backdoors, ASR denotes the proportion of non-empty triggered harmful generations that bypass safety refusal, and lower ASR is better. We use the BackdoorLLM keyword protocol implemented in \texttt{scripts/diagnose\_generation\_metrics.py}: a non-empty triggered output is counted as an attack success if it does not contain any refusal keyword from \texttt{pipeline\_utils.BACKDOORLLM\_JAILBREAK\_KEYWORDS}. This keyword protocol is intentionally auditable, but it has a conservative caveat: incoherent or low-quality non-refusing outputs are still counted as attack successes. We therefore always audit empty outputs and output length together with ASR.

For BEAT, we report HarmRef, BFR, Empty, and output length diagnostics. HarmRef is the harmful-no-trigger refusal rate (also called HNTR in implementation logs): the fraction of non-empty harmful prompts without the trigger that contain a refusal keyword. BFR is the benign-clean false-refusal rate: the fraction of non-empty benign prompts that contain a refusal keyword. Empty is the fraction of generations normalized as empty outputs (\texttt{""}, \texttt{"."}, or \texttt{"?"}); empty outputs are excluded from ASR/HarmRef/BFR denominators and reported separately. Each ASR JSON retains \texttt{total}, \texttt{valid\_total}, \texttt{empty\_count}, \texttt{empty\_output\_rate}, \texttt{avg\_output\_tokens}, and \texttt{median\_output\_tokens}. Higher HarmRef is better, while lower ASR, BFR, Empty, and pathological output shortening are better.

Utility is evaluated with WikiText2 rolling-window token perplexity (PPL) and downstream tasks including BoolQ, RTE, and HellaSwag~\cite{wikitext,boolq,glue,hellaswag}. PPL uses rolling-window token likelihood with max length 1024 and stride 512. We additionally report BEAT's EMD-based detector metrics~\cite{beat}. These metrics measure how separable triggered and clean outputs remain under the detector: AUROC is the area under the receiver-operating-characteristic curve, AP is average precision, and TPR at 5\% FPR is the true-positive rate when the false-positive rate is fixed at 5\%. Lower detector separability after defense indicates that the trigger signature has been weakened.

Runtime, precision, and environment details are reported in Appendix~\ref{app:impl_details} to keep the main text focused on method and experimental findings.

\section{Main Results}

\subsection{BEAT Llama-3.1 Jailbreak Backdoors}
Table~\ref{tab:beat_main} reports the BEAT jailbreak results. Word-trigger ASR decreases from 0.9250 to 0.1417. Phrase-trigger ASR decreases from 0.9000 to 0.0917. Long-trigger ASR decreases to 0.1833 in the fixed base runtime, but residual trigger behavior remains.

\begin{table*}[!t]
\caption{BEAT Llama-3.1 Jailbreak Results}
\label{tab:beat_main}
\centering
\begin{tabular}{lcccl}
\toprule
Trigger & Raw ASR$\downarrow$ & Def. ASR$\downarrow$ & Reduction & Notes \\
\midrule
Word & 0.9250 & 0.1417 & 84.7\% & balanced best \\
Phrase & 0.9000 & \textbf{0.0917} & 89.8\% & paper table run \\
Long & 0.9250 & 0.1833 & 80.2\% & residual signal \\
\bottomrule
\end{tabular}
\end{table*}

\begin{figure}[!t]
\centering
\maybeincludegraphics[width=\linewidth]{fig_asr_by_trigger.pdf}
\caption{Triggered ASR before and after defense on BEAT Llama-3.1 jailbreak backdoors. The grouped bars summarize word, phrase, and long trigger settings under the paper-ready operating points.}
\label{fig:asr_by_trigger}
\end{figure}

\subsection{Utility Preservation}
Table~\ref{tab:utility} reports utility metrics for BEAT models. Under the unified rolling-window PPL protocol, defended checkpoints have higher PPL than the corresponding raw models, indicating a modest language-modeling capacity cost rather than a PPL improvement. Downstream task accuracy remains broadly preserved: RTE improves in these runs, while HellaSwag decreases moderately without catastrophic degradation.

\begin{table*}[!t]
\caption{Utility Metrics on BEAT Models}
\label{tab:utility}
\centering
\begin{tabular}{llcccc}
\toprule
Model & Setting & PPL$\downarrow$ & BoolQ$\uparrow$ & RTE$\uparrow$ & HellaSwag$\uparrow$ \\
\midrule
Word & Raw & \textbf{10.04} & 0.7869 & 0.6354 & 0.5468 \\
Word & Defended & 13.80 & \textbf{0.8031} & \textbf{0.7148} & 0.5166 \\
Phrase & Raw & \textbf{9.59} & 0.7749 & 0.6245 & 0.5511 \\
Phrase & Defended & 13.79 & 0.7648 & \textbf{0.7148} & 0.5122 \\
Long & Raw & \textbf{10.52} & 0.7618 & 0.6101 & 0.5541 \\
Long & Defended & 13.84 & \textbf{0.7752} & \textbf{0.6390} & 0.5101 \\
\bottomrule
\end{tabular}
\end{table*}

\subsection{EMD-Based Detection}
Table~\ref{tab:emd} shows BEAT EMD detector metrics. Word and phrase triggers become close to random distinguishability after recovery, while long triggers retain a stronger output-distribution signal.

\begin{table}[!t]
\caption{BEAT EMD Detection Results}
\label{tab:emd}
\centering
\begin{tabular}{llccc}
\toprule
Trigger & Stage & AUROC & AP & TPR@5\% \\
\midrule
Word & Pruned & 99.87 & 99.74 & 100.0 \\
Word & Recovered & \textbf{50.00} & 33.33 & 0.0 \\
Phrase & Pruned & 99.91 & 99.82 & 100.0 \\
Phrase & Recovered & \textbf{54.32} & 42.00 & 13.0 \\
Long & Pruned & 100.00 & 99.99 & 100.0 \\
Long & Recovered & 86.97 & 87.05 & 76.0 \\
\bottomrule
\end{tabular}
\end{table}

\subsection{Cross-Model and Cross-Mechanism Stress Tests}
Table~\ref{tab:cross_model} summarizes additional stress tests. Mistral-3-7B word-trigger BEAT shows strong transfer after lowering the recovery learning rate and increasing the safety-preservation weight. We also measure Mistral EMD-style detector behavior and utility diagnostics for both Mistral and Trojan-Llama. The Mistral word-trigger setting has the clearest ASR--utility--EMD evidence package, while Mistral phrase/long retain much stronger residual attack behavior. Trojan-Llama-8B shows a measurable but insufficient ASR reduction with mostly preserved utility. We therefore use Mistral word as the main transfer evidence and treat Mistral phrase/long and Trojan-Llama as limitation or future-work cases.

\begin{table*}[!t]
\caption{Cross-Model Stress Tests}
\label{tab:cross_model}
\centering
\begin{tabular}{lccc}
\toprule
Model / Trigger & Raw ASR$\downarrow$ & Best Def. ASR$\downarrow$ & Reduction \\
\midrule
Mistral-3-7B Word & 0.908 & \textbf{0.150} & 83.5\% \\
Mistral-3-7B Phrase & 0.933 & 0.783 & 16.1\% \\
Mistral-3-7B Long & 0.892 & 0.625 & 29.9\% \\
Trojan-Llama-8B & 0.880 & 0.675 & 23.3\% \\
\bottomrule
\end{tabular}
\end{table*}

Table~\ref{tab:mistral_emd} reports Mistral EMD detector metrics, which follow the same trend as Llama: word and phrase detector separability drops markedly after defense, while long retains a stronger signal. Table~\ref{tab:stress_utility} reports utility diagnostics for the cross-model stress tests using the same rolling-window PPL protocol as Table~\ref{tab:utility}. These results show that the weaker cross-model ASR reductions are not caused by universal utility collapse: Mistral long and Trojan-Llama have only small rolling-PPL increases, while Mistral word and phrase show clearer utility costs. We compare PPL only within each model setting; absolute PPL values are not used for cross-family ranking.

\begin{table*}[!t]
\caption{Mistral-3-7B EMD Detection Results}
\label{tab:mistral_emd}
\centering
\begin{tabular}{llcccc}
\toprule
Trigger & Stage & ASR & AUROC & AP & TPR@5\% \\
\midrule
Word & Raw & 0.908 & 99.61 & 99.05 & 100.00 \\
Word & Def & 0.150 & 69.40 & 43.31 & 7.00 \\
Phrase & Raw & 0.933 & 99.69 & 99.19 & 100.00 \\
Phrase & Def & 0.783 & 84.58 & 69.85 & 27.00 \\
Long & Raw & 0.892 & 99.50 & 95.80 & 100.00 \\
Long & Def & 0.625 & 96.66 & 83.74 & 94.00 \\
\bottomrule
\end{tabular}
\end{table*}

\begin{table*}[!t]
\caption{Utility Diagnostics for Cross-Model Stress Tests. PPL uses the rolling-window token protocol. For Mistral-Long, the rolling-PPL row is measured on the reproducible FGSM-proxy checkpoint used in the follow-up diagnostics, so it should be read as a utility-protocol diagnostic rather than a paired measurement of every historical best-ASR checkpoint.}
\label{tab:stress_utility}
\centering
\begin{tabular}{llccccc}
\toprule
Model / Trigger & Stage & PPL$\downarrow$ & BoolQ$\uparrow$ & RTE$\uparrow$ & HellaSwag$\uparrow$ & ASR$\downarrow$ \\
\midrule
Mistral Word & Raw & 10.98 & 0.642 & 0.675 & 0.487 & 0.908 \\
Mistral Word & Def & 17.98 & 0.622 & 0.531 & 0.361 & 0.150 \\
Mistral Phrase & Raw & 10.97 & 0.754 & 0.751 & 0.484 & 0.933 \\
Mistral Phrase & Def & 13.61 & 0.680 & 0.635 & 0.399 & 0.783 \\
Mistral Long & Raw & 12.11 & 0.677 & 0.646 & 0.489 & 0.892 \\
Mistral Long & Def & 12.86 & 0.675 & 0.621 & 0.455 & 0.625 \\
Trojan-Llama & Raw & 8.02 & 0.852 & 0.697 & 0.602 & 0.880 \\
Trojan-Llama & Def & 10.14 & 0.844 & 0.700 & 0.575 & 0.675 \\
\bottomrule
\end{tabular}
\end{table*}

\section{Ablation Studies}

\subsection{Comparison with CROW-Style Direct Repair}
Because our proxy construction is inspired by clean-data perturbation repair, we compare against the public CROW implementation adapted to our BEAT evaluation protocol~\cite{crow}. CROW applies clean-data perturbation and consistency regularization as direct LoRA repair, without structural scoring or pruning. We first observe that the original configuration is unstable on BEAT jailbreak backdoors: low-step training can destroy harmful-prompt refusal, while higher-step training can collapse generation. To avoid an unfairly weak baseline, we further tune learning rate, consistency weight, and training steps on BEAT word, then migrate effective settings to phrase, long, and Mistral stress-test models.

\begin{table*}[!t]
\caption{Tuned CROW Direct Repair vs. Structural Pruning. CROW Effect uses the unified sign convention: positive denotes ASR reduction, negative denotes ASR increase. The ``Ours'' column matches the paper-final main results in Table~\ref{tab:beat_main}.}
\label{tab:crow_baseline}
\centering
\scriptsize
\begin{tabular}{llccccl}
\toprule
Model & Trigger & Raw ASR$\downarrow$ & CROW ASR$\downarrow$ & CROW Effect & Ours ASR$\downarrow$ & Note \\
\midrule
Llama-3.1 & Word & 0.925 & 0.867 & +6.3\% & \textbf{0.1417} & weak CROW reduction \\
Llama-3.1 & Phrase & 0.900 & 0.750 & +16.7\% & \textbf{0.0917} & best CROW case \\
Llama-3.1 & Long & 0.925 & 0.908 & +1.8\% & \textbf{0.1833} & near raw \\
Mistral-3 & Word & 0.908 & 0.950 & -4.6\% & \textbf{0.150} & gen\_len 62$\rightarrow$14 \\
Mistral-3 & Phrase & 0.933 & 0.958 & -2.7\% & \textbf{0.7833} & gen\_len 53$\rightarrow$35 \\
Mistral-3 & Long & 0.892 & 0.967 & -8.4\% & \textbf{0.625} & gen\_len 25$\rightarrow$14 \\
\bottomrule
\end{tabular}
\end{table*}

The complete tuning results show that CROW has a narrow operating window: on Llama word, only learning rate $10^{-4}$ produces a valid effect, while lower learning rates leave the model close to the raw backdoor. The consistency weight has little impact on word-trigger ASR. CROW helps Llama phrase most but remains far weaker than structural pruning, and Llama long is nearly unaffected. On Mistral, migrated CROW direct repair increases ASR and produces clear generation-length collapse. These results suggest that perturbation sensitivity alone is insufficient as a dense repair objective and does not transfer robustly across models; our method instead converts the same class of signal into high-confidence structural scores and removes selected units before masked recovery.

\subsection{Auditable Matched-Random Controls}
To test whether the score identifies useful structural units beyond random pruning and recovery dynamics, we compare the scored pruning plan against auditable matched-random baselines. The reference plan is the high-confidence BEAT word plan selected by the raw score cap $S(u)\le 0$, containing 21 deep MLP channels. We emphasize that this is a focused attribution setting with a tight confidence gate (pruning only about 21 high-confidence units); its purpose is to verify whether scored pruning carries more structural information than matched random controls, so the scored reference value of 0.258 is not the paper-best operating point in Table~\ref{tab:beat_main} (word $=0.1417$) and the two should not be mixed. For each random baseline, we generate independent pruning plans with recorded unit hashes, pairwise Jaccard overlap, layer/component histograms, score statistics, and zero fallback usage. All random plans use the same recovery protocol as the scored plan.

\begin{table*}[!t]
\caption{Auditable Matched-Random Control on BEAT Word. This table uses a focused BEAT-Word attribution setting with tight confidence gating. It evaluates whether scored pruning provides more useful structural information than auditable matched random controls, and is not the paper-best operating point reported in Table~\ref{tab:beat_main}.}
\label{tab:matched_random}
\centering
\begin{tabular}{lccc}
\toprule
Plan & Seeds & Recovered ASR$\downarrow$ & Gap vs. scored \\
\midrule
Scored reference & 1 & \textbf{0.258} & 0.000 \\
Layer/component matched random & 10 & $0.475\pm0.032$ & +0.217 \\
Score-band/component matched random & 10 & $0.516\pm0.027$ & +0.258 \\
Uniform random & 5 & $0.532\pm0.016$ & +0.274 \\
\bottomrule
\end{tabular}
\end{table*}

All 25 recovered random models are worse than the scored reference. In addition, plan-only auditing shows that all 63 random control plans have unique selected-unit hashes, zero fallback usage, and maximum pairwise Jaccard below 0.025. This result shows not only that the scored plan is better for one seed, but also that the proxy score provides useful ranking information beyond structure distribution, score-band distribution, and the recovery protocol. It does not prove exact trigger-unit localization; rather, the more conservative interpretation is that high-confidence perturbation-sensitive units are more useful than structure-matched random units.

Additional follow-up ablations are provided in Appendix~\ref{app:followup_ablations}. These include score normalization, trigger-free checkpoint selection, proxy-safe recovery, recovery alignment strength, and recovery schedule variants. They are summarized in the appendix because they either produce negative or unstable results, or explain the selected operating point without changing the main method.

\section{Discussion}

\subsection{Why the Perturbation Proxy Helps}
The method assumes that some jailbreak backdoor pathways create structural sensitivity to abnormal representation directions. Since the real trigger is unknown, directly searching for trigger samples is infeasible. The adversarial consistency perturbation instead probes the clean-data neighborhood for directions that create large internal shifts. Units that strongly respond to this proxy are treated as high-confidence perturbation-sensitive candidates, not as certified trigger units.

This proxy does not guarantee equivalence to the true trigger. We therefore view the method as mitigation rather than certified removal. Experiments support this position: the proxy is sufficient for strong mitigation on BEAT word and phrase jailbreak settings, and matched-random controls show that the high-confidence scored units are more useful than structurally matched random units. The CROW comparison further shows that directly training on perturbation consistency is not enough; the perturbation signal becomes effective only when converted into confidence-gated structural pruning decisions. However, long-context triggers retain residual signal, and cross-model transfer is partial.

\subsection{Auto-Calibration and Representation Proxies}
The main results in this paper use fixed profiles selected from controlled experiments. For unseen models, fixed profiles may not be reliable. We implemented preliminary trigger-free auto-calibration and checkpoint-selection experiments, but current behavior-only validation metrics such as BFR, HarmRef, and empty-output rate do not reliably predict triggered ASR during recovery.

We also ran follow-up proxy-signal studies on the difficult Mistral-Long setting. Sequence-pooled consistency proxies did not improve over the per-token FGSM proxy, and simple activation-magnitude representation signals selected load-bearing channels and caused collapse. A trigger-aware causal/logit-gap representation signal, however, showed that a separable channel-level circuit exists: using triggered prompts, causal selection plus the same recovery protocol reduced ASR to 0.333 with rolling PPL 12.63. Replacing the real trigger with a trigger-free refusal-suppression perturbation partially closed this gap (best ASR 0.433, PPL 12.97), but at a high benign false-refusal cost (BFR 0.590), and a stronger multi-step PGD soft perturbation did not improve it. Thus the remaining future-work target is not generic sequence pooling, but a balanced trigger-free causal proxy that captures the causal circuit without importing excessive refusal behavior.

\subsection{Limitations}
The method has several limitations. First, it is a mitigation technique rather than a formal guarantee of removal or exact trigger-unit localization. Second, long-context triggers remain difficult. Follow-up experiments indicate that the backdoor can be localized by a trigger-aware causal signal, but the current trigger-free approximation either leaves high ASR or incurs high benign false refusal. Third, BEAT jailbreak results do not directly transfer to all jailbreak injection mechanisms: Mistral phrase/long and Trojan-Llama show only partial response. Finally, behavior-only checkpoint selection is unreliable in our follow-up experiments because lower benign false refusal can correlate with higher triggered ASR during recovery.

\section{Conclusion}
This paper presents Trigger-Free Pruning Defense, a structured mitigation framework for jailbreak-style LLM backdoors under unknown-trigger conditions. The method uses clean benign data and harmful-no-trigger safety prompts, constructs adversarial consistency proxy gradients, scores attention heads and MLP channels, prunes high-confidence perturbation-sensitive structural units, and recovers utility under strict masks. Experiments show that the method reduces BEAT Llama-3.1 word-trigger jailbreak ASR from 0.9250 to 0.1417 and phrase-trigger ASR from 0.9000 to 0.0917, while long-trigger results show meaningful but incomplete mitigation. A tuned CROW-style direct repair baseline remains much weaker, and auditable matched-random controls confirm that the scored high-confidence plan is stronger than structurally matched random pruning. Cross-model stress tests expose important boundaries: not every backdoor mechanism yields an equally useful clean-data proxy signal. Overall, confidence-gated perturbation-sensitive pruning provides an interpretable path for trigger-free jailbreak backdoor mitigation, with clear limits that motivate balanced trigger-free causal or representation-level signals.

\appendices

\section{Implementation Details}
\label{app:impl_details}
BEAT jailbreak experiments use bfloat16 and \texttt{eval\_max\_new\_tokens=64}. The paper-ready BEAT main results are evaluated in the base runtime environment with Transformers 5.3.0. We also observed that the long-trigger result is sensitive to the Transformers runtime: the same checkpoint is approximately 0.18--0.19 ASR under Transformers 5.3.0 and approximately 0.42--0.45 ASR under Transformers 4.57.1. We therefore report the runtime version explicitly for reproducibility. Unless otherwise noted, the reported experiments use Transformers 5.3.0.

Before adding new completeness experiments, we first run a protocol audit and fix the scoring and evaluation protocols for the table being updated. For the BEAT Llama-3.1 Word anchor, the score artifact used for the paper-best 21-unit pruning plan is reproduced with the old \texttt{89d79b1} scoring code, \texttt{chat} prompt template, scoring max length 256, \(\alpha_{\mathrm{safe}}=0.5\), gate \(S(u)\le0\), and \texttt{min\_prune\_layer=2}. The final behavioral evaluation remains the paper protocol: \texttt{alpaca} prompt template, bfloat16, \texttt{eval\_max\_length=1024}, \texttt{eval\_max\_new\_tokens=64}, and greedy decoding. We therefore record score-stage and eval-stage prompt templates separately. New results are not mixed with historical checkpoints produced under a different prompt template or decoding protocol. The audit records dtype, Transformers runtime, decoding settings, and the refusal-keyword source.

For seed stability, the preferred main-table protocol is \texttt{n=3} full-pipeline seeds \(\{13,17,23\}\), rerunning scoring, pruning, recovery, and evaluation for each seed. The seed controls clean/safety prompt sampling, perturbation and scoring randomness, recovery data order, dropout, optimizer randomness, and PyTorch randomness. Evaluation uses greedy decoding and is deterministic under a fixed checkpoint. If compute constraints require a fixed pruning plan, the corresponding table is labeled as mean$\pm$std over recovery seeds with a fixed pruning plan, not as full-pipeline seed stability.

For prune-budget sweeps, the structural-unit denominator is 459,776 for the 32-layer Llama/Mistral 7--8B BEAT models considered here. We sweep requested budgets 0.1\%, 0.3\%, 0.5\%, and 1\%, corresponding to approximately 460, 1379, 2299, and 4598 structural units. Because attention heads and MLP channels have different functional granularity, each sweep row reports requested budget, actual pruned units after gating, number of heads, number of MLP channels, ASR, HarmRef, BFR, Empty, and PPL. Threshold sweeps fix the requested budget at 0.3\% and vary the gate over no gate, \(+0.05\), \(0\), \(-0.02\), and \(-0.05\). Recovery sensitivity is one-factor-at-a-time over steps \(\{10,25,50\}\), learning rate \(\{3\mathrm{e}{-}6,5\mathrm{e}{-}6,1.5\mathrm{e}{-}5,3\mathrm{e}{-}5\}\), \(\lambda_{\mathrm{align}}\in\{1.0,1.5,2.0,2.5\}\), and \(\lambda_{\mathrm{safe}}\in\{0,0.04,0.08,0.12\}\).

We refer to the adaptive robustness check as a Pruning-Aware Reinforcement Stress Test, not as a complete adaptive-attack evaluation. The two variants are \texttt{vanilla\_poisoned\_reinforcement} and \texttt{pruning\_aware\_masked\_reinforcement}. Both start from the released raw BEAT checkpoint and continue reinforcement rather than training a pruning-aware BEAT model from a clean base. The pruning-aware variant masks the top units from the main 0.3\% pruning plan, preferably the actual gate-retained pruned units. If full-budget masking is infeasible, a top-320 mask is reported only as a lower-bound stress test. This experiment tests whether a backdoor can be reinforced to avoid the defender's high-confidence pruning mask; it does not certify robustness against all from-scratch pruning-aware poisoning attacks.

Table~\ref{tab:main_config} lists the recovery configuration for each main result, and Table~\ref{tab:global_config} lists the global settings.

\begin{table*}[!t]
\caption{Recovery Configuration for Each Main Result}
\label{tab:main_config}
\centering
\begin{tabular}{llcl}
\toprule
Model & Trigger & Def. ASR$\downarrow$ & Main config \\
\midrule
Llama-3.1-8B & Word & 0.1417 & $\lambda_{\mathrm{align}}{=}2.0$, $\lambda_{\mathrm{safe}}{=}0.08$, steps${=}25$ \\
Llama-3.1-8B & Phrase & 0.0917 & $\lambda_{\mathrm{safe}}{=}0.06$, ckpt \texttt{beat\_phrase\_ls006}; exact $\lambda_{\mathrm{align}}$/steps not recovered from full log \\
Llama-3.1-8B & Long & 0.1833 & $\lambda_{\mathrm{align}}{=}2.5$, $\lambda_{\mathrm{safe}}{=}0.07$, steps${=}25$ \\
Mistral-3-7B & Word & 0.150 & $\lambda_{\mathrm{safe}}{=}0.20$, lr${=}5\mathrm{e}{-}6$, steps${=}20$ \\
Mistral-3-7B & Phrase & 0.7833 & $\lambda_{\mathrm{safe}}{=}0.15$, lr${=}5\mathrm{e}{-}6$, steps${=}20$ \\
Mistral-3-7B & Long & 0.625 & $\lambda_{\mathrm{safe}}{=}0.30$, lr${=}3\mathrm{e}{-}6$, steps${=}18$ \\
\bottomrule
\end{tabular}
\end{table*}

\begin{table}[!t]
\caption{Global Settings}
\label{tab:global_config}
\centering
\begin{tabular}{ll}
\toprule
Setting & Value \\
\midrule
dtype & bf16 \\
eval max new tokens & 64 \\
Llama main runtime & base, Transformers 5.3.0 \\
PPL protocol & rolling-window token PPL, max length 1024, stride 512 \\
perturbation strength $\epsilon_p$ & 0.1 \\
score weights & $\alpha{=}1.0$, $\beta{=}1.0$ \\
recovery clean weight & $\lambda_{\mathrm{clean}}{=}1.0$ \\
mask policy & strict mask reapplication \\
hardware & 4$\times$ RTX 3090 \\
\bottomrule
\end{tabular}
\end{table}

\section{Additional Follow-up Ablations}
\label{app:followup_ablations}

\subsection{Score Normalization and Recovery Selection}
We tested several lightweight modifications that appeared plausible but did not become part of the main method. Score normalization replaces the raw score with a layer/component-wise $z$-score or rank-normalized variant. Checkpoint selection chooses an intermediate recovery checkpoint using trigger-free validation metrics such as BFR, HarmRef, and empty-output rate instead of always using the final recovery step. Proxy-safe recovery adds a small adversarial safety-preservation term during recovery on harmful prompts without triggers. Table~\ref{tab:negative_followups} reports the ASR outcomes. Note that the raw column in Table~\ref{tab:negative_followups} corresponds to a focused follow-up operating point whose absolute ASR is higher than the paper-best result in Table~\ref{tab:beat_main} (e.g., word 0.317, phrase 0.292, long 0.425); the two should not be mixed.

\begin{table*}[!t]
\caption{Negative or Unstable Follow-up Ablations}
\label{tab:negative_followups}
\centering
\begin{tabular}{llcc}
\toprule
Experiment & Setting & Baseline / Final ASR$\downarrow$ & Variant ASR$\downarrow$ \\
\midrule
Score normalization & Word raw / z / rank & \textbf{0.317} & 0.508 / 0.525 \\
Score normalization & Phrase raw / z / rank & \textbf{0.292} & 0.325 / 0.333 \\
Score normalization & Long raw / z / rank & \textbf{0.425} & 0.450 / 0.425 \\
Checkpoint selection & Phrase final vs selected & \textbf{0.167} & 0.308 \\
Checkpoint selection & Long final vs selected & \textbf{0.258} & 0.383 \\
Proxy-safe recovery & Long baseline vs best & 0.425 & \textbf{0.375} \\
\bottomrule
\end{tabular}
\end{table*}

Raw scores outperform threshold-matched normalized scores on all three BEAT triggers. The reason is semantic: $S(u)\le0$ is a high-confidence boundary for the calibrated raw score, but the same threshold corresponds roughly to a central score region under $z$-score or rank normalization and therefore selects too many units. Trigger-free checkpoint selection also fails to predict triggered ASR during recovery: it tends to select early checkpoints with lower BFR but higher triggered ASR. Proxy-safe recovery gives a small improvement on Long at a small perturbation radius, but the effect is unstable and larger $\epsilon$ values undo the defense. We further tested the optional harmful-context proxy term; it gives a small pruned-only improvement for the phrase trigger, but is inconsistent on word and long triggers and does not reliably survive recovery. For these reasons, harmful-context proxy scoring, proxy-safe recovery, and checkpoint selection are kept out of the main method.

\subsection{Representation and Trigger-Free Causal Follow-ups}
Table~\ref{tab:causal_followup} summarizes the Mistral-Long representation follow-ups. The trigger-aware causal signal is not deployable under our threat model because it uses triggered prompts, but it is a useful upper-bound diagnostic: it shows that a channel-level long-trigger circuit can be isolated without the collapse caused by activation-magnitude signals. The best trigger-free causal approximation improves ASR substantially over the FGSM baseline but is not balanced enough for the main method because BFR rises to 0.590. Multi-step PGD perturbations reduce this benefit.

\begin{table*}[!t]
\caption{Mistral-Long Representation and Trigger-Free Causal Follow-ups}
\label{tab:causal_followup}
\centering
\begin{tabular}{lcccc}
\toprule
Variant & ASR$\downarrow$ & HarmRef$\uparrow$ & BFR$\downarrow$ & Rolling PPL$\downarrow$ \\
\midrule
Raw & 0.892 & 0.400 & 0.160 & 12.11 \\
FGSM proxy baseline & 0.683 & 0.467 & 0.220 & 12.86 \\
Trigger-aware causal K=512 & 0.450 & 0.433 & 0.300 & 12.50 \\
Trigger-aware causal K=1024 & \textbf{0.333} & 0.450 & 0.420 & 12.63 \\
Trigger-free refsupp $\epsilon=0.01$, K=1024 & 0.433 & \textbf{0.667} & 0.590 & 12.97 \\
Trigger-free PGD $\epsilon=0.03$, K=512 & 0.642 & 0.383 & \textbf{0.200} & 13.79 \\
\bottomrule
\end{tabular}
\end{table*}

\subsection{Recovery Alignment Strength}
Table~\ref{tab:lambda_ablation} shows a nonlinear transition in BEAT word experiments. Increasing $\lambda_{\mathrm{align}}$ from 1.0 to 1.5 worsens ASR, while 2.0 sharply improves ASR. This suggests that recovery has a phase-transition-like regime: insufficient proxy alignment allows backdoor features to re-emerge, while sufficiently strong alignment stabilizes internal representations.

\begin{table}[!t]
\caption{BEAT Word $\lambda_{\mathrm{align}}$ Ablation}
\label{tab:lambda_ablation}
\centering
\begin{tabular}{cccc}
\toprule
$\lambda_{\mathrm{align}}$ & ASR$\downarrow$ & HarmRef$\uparrow$ & BFR$\downarrow$ \\
\midrule
1.0 & 0.925 & 0.500 & 0.100 \\
1.5 & 0.950 & 0.450 & 0.100 \\
2.0 & \textbf{0.142} & 0.875 & 0.100 \\
\bottomrule
\end{tabular}
\end{table}

\begin{figure}[!t]
\centering
\maybeincludegraphics[width=\linewidth]{fig_lambda_align_phase.pdf}
\caption{Phase-transition-like behavior of recovery alignment strength on BEAT word jailbreak mitigation. Moderate alignment can be unstable, while sufficiently strong alignment suppresses backdoor relapse during masked recovery.}
\label{fig:lambda_align_phase}
\end{figure}

\subsection{Objective Schedule}
Table~\ref{tab:schedule_ablation} compares recovery schedules using reproducible rerun results. Simultaneous optimization with $\lambda_{\mathrm{align}}=2.0$ gives the strongest balanced point among the evaluated schedules, while alternating variants reduce benign false refusal in some cases but leave higher triggered ASR.

\begin{table}[!t]
\caption{Objective Schedule Ablation on BEAT Word}
\label{tab:schedule_ablation}
\centering
\begin{tabular}{lccc}
\toprule
Schedule & ASR$\downarrow$ & HarmRef$\uparrow$ & BFR$\downarrow$ \\
\midrule
Simultaneous & \textbf{0.142} & 0.875 & 0.100 \\
Alternating & 0.325 & 0.775 & \textbf{0.060} \\
Alternating-soft & 0.258 & 0.825 & 0.080 \\
\bottomrule
\end{tabular}
\end{table}

\begin{figure}[!t]
\centering
\maybeincludegraphics[width=\linewidth]{fig_schedule_tradeoff.pdf}
\caption{ASR--BFR trade-off under different recovery objective schedules. The plot highlights the practical operating point used in the main result and compares only reproducible rerun results.}
\label{fig:schedule_tradeoff}
\end{figure}

\begin{thebibliography}{00}

\bibitem{badnets}
T. Gu, B. Dolan-Gavitt, and S. Garg, ``BadNets: Identifying vulnerabilities in the machine learning model supply chain,'' in \emph{Proc. Workshop on Machine Learning and Computer Security}, 2017.

\bibitem{trojaning}
Y. Liu, S. Ma, Y. Aafer, W.-C. Lee, J. Zhai, W. Wang, and X. Zhang, ``Trojaning attack on neural networks,'' in \emph{Proc. Network and Distributed System Security Symposium (NDSS)}, 2018.

\bibitem{hidden_trigger}
A. Saha, A. Subramanya, and H. Pirsiavash, ``Hidden trigger backdoor attacks,'' in \emph{Proc. AAAI Conference on Artificial Intelligence}, 2020.

\bibitem{universal_triggers}
E. Wallace, S. Feng, N. Kandpal, M. Gardner, and S. Singh, ``Universal adversarial triggers for attacking and analyzing NLP,'' in \emph{Proc. EMNLP-IJCNLP}, 2019.

\bibitem{autoprompt}
T. Shin, Y. Razeghi, R. L. Logan IV, E. Wallace, and S. Singh, ``AutoPrompt: Eliciting knowledge from language models with automatically generated prompts,'' in \emph{Proc. EMNLP}, 2020.

\bibitem{zou2023universal}
A. Zou, Z. Wang, J. Z. Kolter, and M. Fredrikson, ``Universal and transferable adversarial attacks on aligned language models,'' arXiv:2307.15043, 2023.

\bibitem{gptfuzzer}
Y. Yu, X. Lin, Z. Yu, and X. Xing, ``GPTFuzzer: Red teaming large language models with auto-generated jailbreak prompts,'' arXiv:2309.10253, 2023.

\bibitem{jailbroken}
A. Wei, N. Haghtalab, and J. Steinhardt, ``Jailbroken: How does LLM safety training fail?,'' in \emph{Proc. NeurIPS}, 2023.

\bibitem{sleeper_agents}
E. Hubinger \emph{et al.}, ``Sleeper agents: Training deceptive LLMs that persist through safety training,'' arXiv:2401.05566, 2024.

\bibitem{neuralcleanse}
B. Wang, Y. Yao, S. Shan, H. Li, B. Viswanath, H. Zheng, and B. Y. Zhao, ``Neural Cleanse: Identifying and mitigating backdoor attacks in neural networks,'' in \emph{Proc. IEEE Symposium on Security and Privacy}, 2019.

\bibitem{spectral}
B. Tran, J. Li, and A. Madry, ``Spectral signatures in backdoor attacks,'' in \emph{Proc. NeurIPS}, 2018.

\bibitem{activation_clustering}
B. Chen, W. Carvalho, N. Baracaldo, H. Ludwig, B. Edwards, T. Lee, I. Molloy, and B. Srivastava, ``Detecting backdoor attacks on deep neural networks by activation clustering,'' arXiv:1811.03728, 2018.

\bibitem{finepruning}
K. Liu, B. Dolan-Gavitt, and S. Garg, ``Fine-pruning: Defending against backdooring attacks on deep neural networks,'' in \emph{Proc. RAID}, 2018.

\bibitem{fgsm}
I. J. Goodfellow, J. Shlens, and C. Szegedy, ``Explaining and harnessing adversarial examples,'' in \emph{Proc. ICLR}, 2015.

\bibitem{pgd}
A. Madry, A. Makelov, L. Schmidt, D. Tsipras, and A. Vladu, ``Towards deep learning models resistant to adversarial attacks,'' in \emph{Proc. ICLR}, 2018.

\bibitem{adv_training_free}
A. Shafahi \emph{et al.}, ``Adversarial training for free!,'' in \emph{Proc. NeurIPS}, 2019.

\bibitem{magnitude_pruning}
S. Han, J. Pool, J. Tran, and W. Dally, ``Learning both weights and connections for efficient neural networks,'' in \emph{Proc. NeurIPS}, 2015.

\bibitem{lottery}
J. Frankle and M. Carbin, ``The lottery ticket hypothesis: Finding sparse, trainable neural networks,'' in \emph{Proc. ICLR}, 2019.

\bibitem{head_pruning}
P. Michel, O. Levy, and G. Neubig, ``Are sixteen heads really better than one?,'' in \emph{Proc. NeurIPS}, 2019.

\bibitem{movement_pruning}
V. Sanh, T. Wolf, and A. M. Rush, ``Movement pruning: Adaptive sparsity by fine-tuning,'' in \emph{Proc. NeurIPS}, 2020.

\bibitem{llm_pruner}
M. Ma, G. Fang, X. Wang, and X. Wang, ``LLM-Pruner: On structural pruning of large language models,'' in \emph{Proc. NeurIPS}, 2023.

\bibitem{wikitext}
S. Merity, C. Xiong, J. Bradbury, and R. Socher, ``Pointer sentinel mixture models,'' in \emph{Proc. ICLR}, 2017.

\bibitem{boolq}
C. Clark, K. Lee, M.-W. Chang, T. Kwiatkowski, M. Collins, and K. Toutanova, ``BoolQ: Exploring the surprising difficulty of natural yes/no questions,'' in \emph{Proc. NAACL-HLT}, 2019.

\bibitem{glue}
A. Wang \emph{et al.}, ``GLUE: A multi-task benchmark and analysis platform for natural language understanding,'' in \emph{Proc. ICLR}, 2019.

\bibitem{hellaswag}
R. Zellers, A. Holtzman, Y. Bisk, A. Farhadi, and Y. Choi, ``HellaSwag: Can a machine really finish your sentence?,'' in \emph{Proc. ACL}, 2019.

\bibitem{beat}
B. Yi, T. Huang, S. Chen, T. Li, Z. Liu, Z. Chu, and Y. Li, ``Probe before you talk: Towards black-box defense against backdoor unalignment for large language models,'' in \emph{Proc. ICLR}, 2025.

\bibitem{crow}
N. M. Min, L. H. Pham, Y. Li, and J. Sun, ``CROW: Eliminating backdoors from large language models via internal consistency regularization,'' in \emph{Proc. ICML}, 2025.

\end{thebibliography}

\end{document}
